\section{Background \& Preliminaries}
\label{sec:background}

\subsection{The Serial Tool-Wait Bottleneck}
\label{sec:bottleneck}

The dominant execution model for LLM agents follows the ReAct
loop~\citep{yao_react_2023}: the model reasons (Thought), emits a structured tool
invocation (Action), waits for the external service to return a result (Observation),
and then resumes reasoning conditioned on that result.
This Thought--Action--Observation cycle enforces a strict serial dependency:
\[
[\text{(Re)prefill}] \;\to\; [\text{Decode}] \;\to\; [\text{Tool wait}] \;\to\; \cdots
\]
Each turn (re)prefills the context (the user prompt on the first turn, the appended tool
result thereafter), decodes the reasoning and tool call, then stalls on the tool wait.
The tool call cannot be dispatched until decoding completes; the next decode cannot begin
until the tool result arrives and is (re)prefilled into the context window.
This creates a structural pipeline stall regardless of how fast individual components are.

\parabf{Tool latency dominates agent wall time.}
Figure~\ref{fig:ttool_dist} shows that tool execution consumes a substantial share of
end-to-end agent latency across three workloads: 16\% on tau2-bench (enterprise database
APIs with 2\,s simulated floor), 19\% on GAIA (real web search APIs), and 37\% on
BrowseComp (web search + page browse).
On BrowseComp, the distribution is heavily right-skewed: median tool latency is 1.19\,s
but $P_{95}$ reaches 70\,s, meaning tail-latency optimization requires addressing tool
stalls, not just model serving.

\parabf{Why existing solutions fall short.}
Token-level speculative decoding~\citep{leviathan_fast_2023,li_eagle_2024,cai_medusa_2024}
reduces generation latency but does not dispatch external tools earlier.
Parallel tool calling (e.g., OpenAI function calls) helps only when the model can identify
independent tools upfront; empirically, only 15--25\% of tool calls in GAIA and SWE-bench~\citep{jimenez_swebench_2024}
are parallelizable with their predecessor.
Workflow optimizers~\citep{orion,stojkovic2023specfaas,fu2024serverlessllm} prewarm known
execution graphs, but agent tool calls are generated \emph{online}: there is no static
graph to prefetch.
Recent agent-speculation systems~\citep{ye_speculative_2026,zhong_dualspec_2026,sui_act_2026}
move closer to execution-level overlap but require separate predictors, verifier policies,
or historical traces.
None speculates from the currently running model alone.

\subsection{Cost Model for Speculative Overlap}
\label{sec:eq1}

If we could predict the next tool call before the model finishes generating it, we could
dispatch the tool speculatively and overlap its execution with ongoing model decode.
The realized speedup depends on four measurable parameters:
\begin{itemize}
  \item $\alpha$: the fraction of tool-call \emph{turns} where speculation is accepted
        (gate rate). $\alpha$ is always measured per turn, not per dispatched probe; a
        retrying gate may issue several probes within one turn.
  \item $\toverlap$: mean realized time overlap per accepted turn (how much tool latency
        is hidden behind decode).
  \item $\Tbasestar$: the realized main-stream decode cost on the speculative path. It
        equals the baseline $\Tbase$ unless prefix recovery is active, which lowers it by
        letting the main stream skip re-decoding the verified tool-call tokens it recovers
        from the rejected probe on a miss.
  \item $\Toh$: per-turn overhead from probe computation and wasted speculative execution
        on rejected turns.
\end{itemize}

\noindent The resulting latency-ratio formula (full derivation in Appendix~\ref{app:eq1}):
\begin{equation}
  \boxed{\mathrm{Ratio} = \frac{\Tbase}{\Tbasestar - \alpha\cdot\toverlap + \Toh}}
  \label{eq:eq1}
\end{equation}

\noindent Speculation improves latency when $\mathrm{Ratio} \geq 1$, which approximately
requires $\alpha\cdot\toverlap \geq \Toh$ (exact form in Appendix~\ref{app:eq1}): the time
saved by overlapping tool execution on accepted turns must exceed the overhead incurred
across all turns.
The upper bound is set by the tool-time fraction $\ftool$: with perfect prediction
($\alpha=1$, $\Toh=0$), $\mathrm{Ratio} = 1/(1-\ftool)$.
For BrowseComp ($\ftool=0.366$), this ceiling is $1.58\times$.

The next section establishes empirically that these parameters are achievable: $\alpha$
is high, $\Toh$ is low, and even rejected probes contribute via prefix recovery.
